% answers_short.tex —— 2012 年·北京卷·理科 · 简短答案速查
% 由 tools/make_short_answers.py 自 answers.tex 的 \gksolhead/\gkanswer 提取，
% 并与 2012/2012_数学真题_北京理_答案要点.md 逐题比对后生成，勿手改（改源 MD 后重跑）。

\begin{gknotice}
\textbf{使用说明}：本节只给最终答案便于快速核对，完整推导见解析版。答案为编写者独立推导，\textbf{非官方标准答案}。
\end{gknotice}

\gkpart{选择题}
\begin{center}\begin{tabular}{ccccccccc}
\toprule
题号 & 1 & 2 & 3 & 4 & 5 & 6 & 7 & 8 \\
\midrule
答案 & D & D & B & C & A & B & B & C \\
\bottomrule
\end{tabular}\end{center}

\gkpart{填空题}
\begin{itemize}[leftmargin=2.6em]
  \item[\textbf{9.}] $2$。
  \item[\textbf{10.}] $a_2=1$；$S_n=\dfrac{n(n+1)}4$。
  \item[\textbf{11.}] $4$。
  \item[\textbf{12.}] $\sqrt3$。
  \item[\textbf{13.}] $1$；$1$。
  \item[\textbf{14.}] $(-4,-2)$。
\end{itemize}

\gkpart{解答题}
\begin{itemize}[leftmargin=2.6em]
  \item[\textbf{15.}] (1) 定义域 $\{x\mid x\neq k\pi,\ k\in\mathbf Z\}$，最小正周期 $T=\pi$；(2) $\left[-\dfrac\pi8+k\pi,\ k\pi\right)$ 与 $\left(k\pi,\ \dfrac{3\pi}8+k\pi\right]$（$k\in\mathbf Z$）
  \item[\textbf{16.}] (1) 证明见解析；(2) $45^\circ$（即 $\dfrac\pi4$）；(3) 不存在
  \item[\textbf{17.}] (1) $\dfrac23$；(2) $\dfrac3{10}$；(3) $a=600$，$b=0$，$c=0$；$s^2=80000$
  \item[\textbf{18.}] (1) $a=3$，$b=3$；(2) 递增区间 $\left(-\infty,-\dfrac a2\right)$ 与 $\left(-\dfrac a6,+\infty\right)$，递减区间 $\left(-\dfrac a2,-\dfrac a6\right)$；最大值当 $0<a<2$ 时为 $a-\dfrac{a^2}4$，当 $a\geqslant2$ 时为 $1$
  \item[\textbf{19.}] (1) $\left(\dfrac72,5\right)$；(2) 证明见解析
  \item[\textbf{20.}] (1) $0.7$；(2) $1$；(3) $\dfrac{2t+1}{t+2}$
\end{itemize}
